% Converted from S5_Tool_call_audit/reviewer_response_package_20261119_claude_v13/RESPONSE_R2_16_EN.md (2026-11-17)
We thank the reviewer for this comment and agree that the original manuscript did not evaluate failure modes of the planner--executor workflow. We have therefore performed a fully instrumented failure-mode audit (\responsetableref{tab_rep_failure_modes} and \responsefigureref{fig_rep_execution_traces}) and summarize its main results in Results~\SecResReport{}.\par

\textbf{Audit design.} We audited 150 report-generation cases, 50 from each of SMU-HMC, KMVE and ZJH-TS. The cases were drawn from a sample stratified by the number of thyroid nodules or cysts stated in the reference report: 0, 1, $\geq$2 or uncertain in 11, 12, 23 and 4 cases in SMU-HMC, 12, 11, 23 and 4 in KMVE, and 1, 11, 38 and 0 in ZJH-TS. Each case was run through the complete workflow three times, giving 450 runs and 2,347 tool calls. We used a root seed of 42 to deterministically derive a pre-assigned seed for each run. For every run we recorded the LLM requests and responses, each tool call with its arguments and raw return, the summary returned to the LLM, the case state before and after each call, and the final report. Each failure category was defined operationally and assessed for all tool calls (\responsetableref{tab_rep_failure_modes}). Because each case was run three times, we also report how many cases were affected and whether an event recurred across a case's runs (\responsetableref{tab_rep_failure_modes}). Confidence intervals (CIs) were obtained by resampling cases within each dataset, keeping all three runs of each case.\par

\textbf{Frequency of the requested failure categories} (\responsetableref{tab_rep_failure_modes}).\par

\begin{itemize}[leftmargin=1.5em,itemsep=1pt,topsep=2pt]
\item \textbf{Invalid calls.} Four of the 2,347 tool calls (0.17\%; 95\% CI 0.04--0.35\%) were invalid. They occurred in four different cases, each in only one of its three runs. One call used a malformed tool name. The other three were gland-measurement requests that cited image identifiers that do not exist: one was a zero-padded variant of an existing identifier, and two were numbered one beyond the last image. The executor did not process the non-existent identifiers.
\item \textbf{Unnecessary calls.} We counted requests that listed images outside the tool's target region: 18 lymph-node screening calls that included non-neck images (0.77\% of all calls; 95\% CI 0.27--1.37\%; 8 cases) and 66 gland-measurement calls that included neck images (2.81\%; 1.96--3.70\%; 32 cases). In both situations the executor passes only the eligible images to the tool, so the other images are simply not analysed in that call; nothing is deleted. For example, one screening request listed ten images, and only its five neck images were screened (\responsefigureref{fig_rep_execution_traces}e). Because each of these calls still analysed eligible images, we regard them as imprecise requests rather than unnecessary calls. Region labels come from automated preprocessing.
\item \textbf{Hallucinated tasks.} We defined a hallucinated task as a task that was planned or executed without support in the case's images. Every one of the 2,347 tool calls executed a task in the plan. We then checked each of the 2,348 planned tasks against the images available at planning time; for example, nodule segmentation requires an image flagged by preprocessing, and lymph-node screening requires a neck image. Three planned tasks (0.13\%; 95\% CI 0.00--0.29\%) had no such support. All three were lymph-node screening planned for cases without neck images, each in a different case and in only one of its three runs. In two of these runs the task was never executed; in the third it was issued with an empty image list and returned nothing.
\item \textbf{Repeated calls.} Across the 450 runs, apart from one retry triggered by a malformed tool name, no tool was called more than once within the same run.
\item \textbf{Execution failures.} The only failed tool call was the malformed name. All 450 runs were completed. Three of them, in three different cases, were first interrupted by out-of-memory errors caused by an external factor (other jobs occupying GPU memory) and were re-run with the same inputs.
\item \textbf{Retries and recovery.} The call with the malformed tool name was retried by the LLM at the next step, with the corrected name and the same arguments. In addition, the executor has a fallback for steps at which the LLM does not continue with the plan produced by the planner, that is, returns neither a tool call nor the completion signal: it then runs the next planned task that has not yet been performed. This fallback was triggered in two runs, each from a different case. In both, the LLM returned an empty response after gland measurement had completed; the executor ran the planned lymph-node screening, and the LLM then completed the run. Because the fallback can only run tasks listed in the plan, it cannot restore a step that the plan itself omits (see the omitted fusion step below).
\item \textbf{Rollback.} The current system does not implement transactional rollback of tool actions.
\end{itemize}

\textbf{Major failure modes and representative traces} (\responsetableref{tab_rep_failure_modes}; \responsefigureref{fig_rep_execution_traces}). Apart from the request errors above, all of which were contained, the main failure modes were ordering errors and omitted steps that the run-level success flags did not reveal:\par

\begin{itemize}[leftmargin=1.5em,itemsep=1pt,topsep=2pt]
\item In 40 of the 320 fusion-tool calls (24 cases; in all three runs of 3 of them), fusion was executed before pixel spacing was available for the nodule-positive images, and spacing obtained later did not trigger re-fusion. In 4 SMU-HMC and ZJH-TS runs, this left 5 nodule descriptions without a size. The other 36 calls were in KMVE cases, whose reports use a mask in place of specific nodule sizes. In addition, in one run, fusion was executed before nodule feature classification.
\item For the 43 cases in which preprocessing flagged no nodule, the workflow follows a plan that does not include spacing estimation. Consequently, gland measurement in these cases was performed without a pixel-spacing estimate.
\item In one run, a segmented nodule was never fused because fusion was missing from the plan, and the report contained no description of that nodule.
\end{itemize}

\responsefigureref{fig_rep_execution_traces} shows real traces for a case without nodules, a case with multiple nodules, conflicting intermediate evidence (a preprocessing-flagged image not confirmed by segmentation and excluded downstream, seen in all three runs of two cases), recovery from a malformed tool name, off-target images that were not processed, fusion before pixel spacing, and the omitted fusion step.\par
